\section{Evaluation}
\subsection{Questions and methodology}
The evaluation asks four questions. Does separating the control boundaries
improve completion time against a modern data-aware baseline? How much of any
improvement comes from fixed concurrency rather than the architecture? Does
deferred preparation change where tasks wait? Do correctness and recovery
remain intact when workers disappear?

The primary artifact contains two campaigns. A colocated campaign requests an
eight-core compute allocation and runs two four-core workers within it. A
separate campaign obtains two or four four-core workers through the batch
system, with the service on a submission host. Worker count is not automatically
physical node count: placement may colocate workers, and every result records
the actual execution hosts. Consequently, the latter is a worker-pool study
on an opportunistic cluster, not an exclusive-node strong-scaling claim.
The colocated cold-library allocation uses an AMD EPYC 7543 host; the
separate preloaded allocation uses an AMD EPYC 9334 host. Python is 3.10.20
and NumPy 1.26.4. Runtime comparisons are within their respective allocation;
cross-campaign differences also include this hardware change.
Each worker is explicitly pinned to exactly its declared number of logical
CPUs; local workers use disjoint CPU sets, and the colocated campaign process
is limited to eight CPUs overall. This is necessary because the batch system
assigns CPU shares without a hard quota. Task records verify actual affinity.
Affinity limits concurrency capacity but does not establish exclusive cores;
other allocations can still contend on the same host.

Every trial creates a fresh service or manager and admits the complete
requested worker pool before starting the measured interval. Timing includes
graph construction and submission, function execution, fetching every sink,
and a local fsync of each fetched sink. DataVine additionally performs its
controller's requested-output durable admission. This is a conservative extra
cost for DataVine, but it means the two implementations do not perform identical
numbers of persistence operations. Intermediate backup is disabled in both
primary comparisons; default-background-backup behavior is tested separately.
Pool admission time is reported separately and excluded from performance.

Within each workload and worker count, variant order is randomized in three
repetition blocks using a retained seed. A trial is accepted only if its
scientific invariants, logical identities, sink checksums, and DataVine physical
attempt counts pass. Failed and interrupted attempts remain in the artifact.
An incomplete pool never becomes a performance sample, and changing executable
inputs prevents resuming a campaign under its old identity. Figures use only
validated runs and expose the number of repetitions. We show ranges with
medians rather than claiming precise population confidence from three samples.

\subsection{Workloads and controls}
Each graph contains 32 lanes and four stages, for 128 distinct task identities.
Every task produces 256\,KiB of output, except where a numerical representation
requires its corresponding element count. Both runtimes invoke the same
serialized Python function. Parent payload hashes are checked at each consumer,
and the sink closure must contain every expected logical identity. Telemetry
records travel with results in both implementations; this is a measurable
instrumentation cost and grows with graph depth.

\begin{table}[t]
\centering\small
\caption{Scientific kernels and controlled stressor. These are miniapplications,
not full production scientific analyses.}
\label{tab:workloads}
\begin{tabular}{@{}p{.18\columnwidth}p{.39\columnwidth}p{.31\columnwidth}@{}}
\toprule
Workload & Computation and graph & Independent check \\
\midrule
Spectral & FFT low-pass pipelines & Analytical signal energy \\
Histogram & Byte scans; adjacent-lane fan-in & Count conservation and hashes \\
Quadrature & Simpson integration pipelines & Analytic Gaussian integral \\
Phase & CPU / disk / CPU / disk, with fan-in & All identities and content hashes \\
\bottomrule
\end{tabular}
\end{table}

Spectral filtering uses NumPy FFTs on a deterministic two-frequency signal;
energy has an analytical expected value. Quadrature computes a smooth Gaussian
integral with composite Simpson integration and checks against the error-function
solution. Histogram tasks scan all parent bytes and check count conservation.
The phase control alternates a 100\,ms process-CPU loop with actual node-local
file writes, fsync, and reads. It does not replace I/O with sleeping, although
its files can still benefit from kernel caching and storage-device behavior.
Its purpose is to expose a resource-phase transition, not model an application
filesystem quantitatively.

The colocated campaign compares DataVine fixed windows of $C$, $2C$, and $4C$,
elastic admission, and elastic admission with eager preparation. TaskVine is
measured at $C$ and with the explicitly extended fixed $2C$/$4C$ controls.
The worker-pool campaign uses four representative variants: DataVine $C$ and
elastic, stock TaskVine $C$, and extended TaskVine $4C$. These comparisons
separate application correctness from scheduling and preparation choices.

\subsection{Completion time and concurrency attribution}
\input{results/evaluation-summary.tex}
\begin{figure*}[t]
\centering\includegraphics[width=.96\textwidth]{completion.pdf}
\caption{Colocated compute-allocation trials. Points show individual accepted
runs; bars indicate medians and whiskers the observed range. TV+ denotes the
private fixed-overcommit extension, not stock TaskVine. Lower is better.}
\label{fig:completion}
\end{figure*}

A comparison with stock TaskVine alone cannot distinguish more available
function slots from cheaper data coordination. The fixed-window controls
address that ambiguity. For each workload, the most informative pair is
DataVine elastic versus the best tested fixed setting within DataVine, followed
by the best tested extended TaskVine setting. If elastic matches a fixed
setting, the result supports convenience across phases only to the extent
that one fixed setting does not already work well everywhere. A negative
result is retained rather than explained away by choosing a different baseline.

The comparison is still not a complete factorial decomposition of scheduler
and controller implementation costs. Both runtimes differ in APIs, graph
submission, serialization, and result admission. Attempt timelines isolate
preparation ordering within DataVine, but an architecture-level causal claim
needs additional matched instrumentation of both control paths. We therefore
report end-to-end differences separately from mechanism observations.

\subsection{Separating dependency deployment from runtime behavior}
The fresh-library campaign includes importing dependencies in each forked
function process. Process-state diagnostics found a substantial deployment
confound: Python children blocked in NFS open and RPC waits. Four additional
identically configured quadrature diagnostics completed in 32.453, 0.711,
0.694, and 1.367 s. The slow run included repeated kernel wait states associated
with NFS; these samples establish an observed source of delay, not an
attribution of every campaign outlier to the filesystem. The function body's
CPU timer begins after its initial standard-library imports and therefore
cannot explain the entire task residence interval.

We consequently add a separate, symmetric-preload control. Both execution
libraries import the same explicit standard-library modules and NumPy before
reporting readiness and before forking function calls. TaskVine uses its
existing module-hoisting interface; DataVine uses a private executor selected
through its existing executor-path override. All later task bodies, data
checks, CPU limits, graph shapes, and persistence policies are unchanged.
The timer follows full worker connection, without a separate all-libraries-ready
barrier; residual library setup can therefore enter the measured interval.
This measures an explicitly configured preload deployment; the original cold-library observations
remain visible and are not replaced or pooled with these runs. The preloaded
campaign uses a separate opportunistic allocation, so ratios between cold and
preloaded campaigns are not paired causal estimates.

\begin{figure*}[t]
\centering\includegraphics[width=.96\textwidth]{preloaded.pdf}
\caption{Matched explicit dependency preloading, with three randomized blocks.
Points show all accepted runs, bars medians, and whiskers observed ranges.
Both backends preload dependencies once per library; residual startup is
included if it outlasts worker admission.}
\label{fig:preloaded}
\end{figure*}
\input{results/preloaded-summary.tex}
These comparisons test the runtime under a controlled library environment.
They do not establish that feedback is responsible for all differences:
DataVine fixed $C$ is the direct counterexample when it matches or exceeds
elastic admission. A production deployment must also account for preparing
and distributing these libraries, and for libraries that are unsafe to fork.

\subsection{Worker-pool sensitivity}
\begin{figure}[t]
\centering\includegraphics[width=\columnwidth]{worker-pool.pdf}
\caption{Completion time versus admitted worker count on the opportunistic
cluster. Host placement is recorded per trial. This is not an exclusive-node
scaling result. Markers are medians; whiskers are observed ranges.}
\label{fig:pool}
\end{figure}
Increasing the worker count can reduce available work per worker while adding
control and transfer endpoints. At 128 tasks, this study tests sensitivity to
pool size rather than maximum system capacity. It is especially important not
to infer scalability from the mere ability to connect more workers. Changes
in physical placement, node-local I/O, and submission-host load can affect
observations even after full admission. The artifact retains those conditions
so follow-up reserved-allocation runs can test the same hypotheses.

\subsection{Data paths, larger graphs, and resource pressure}
\begin{figure}[t]
\centering\includegraphics[width=\columnwidth]{data-pipeline.pdf}
\caption{Larger hash-data pipeline in two separate CPU-limited allocations,
with fresh libraries and with explicit dependency preloading.
All points represent validated complete runs; TV+ is the fixed-overcommit
extension. DV E and DV P denote elastic and eager preparation. The kernel reads,
hashes, and generates payloads.}
\label{fig:data-pipeline}
\end{figure}
\begin{figure}[t]
\centering\includegraphics[width=\columnwidth]{memory-pressure.pdf}
\caption{Independent memory-pressure study: completion time and peak sampled
worker memory for a declared 128-MiB task footprint. Sampling is not a hard
bound on every allocation between observations.}
\label{fig:memory}
\end{figure}
\begin{figure}[t]
\centering\includegraphics[width=\columnwidth]{routes.pdf}
\caption{Source-isolated delivery of 128 files. Admission and backup preparation
precede the measured consumer phase. Network utilization is not inferred from
shared frontend counters.}
\label{fig:routes}
\end{figure}
\input{results/supplemental-summary.tex}

\subsection{Preparation timeline and exercised failures}
\begin{figure}[t]
\centering\includegraphics[width=\columnwidth]{preparation.pdf}
\caption{Worker-local residence decomposition for the focused 32-task
functional experiment, using median per-attempt intervals; component medians need not sum to the
median total. This submission-host
run validates and illustrates the instrumentation; it is not a distributed
performance result.}
\label{fig:preparation}
\end{figure}
The focused eager/deferred experiment verifies 32 ordered traces in each mode.
It shows where residence can be measured without conflating pre-preparation
queueing with post-preparation waiting. These short, single-worker runs do not
establish a throughput advantage for either policy. In particular, delaying
preparation may sacrifice useful prefetch when transfers are long; the main
campaign's eager variant tests the end-to-end effect independently.

\input{results/recovery-summary.tex}
The regression suite also exercises protocol, lifecycle, and persistence
boundaries. Passing these cases supports the stated tested semantics; it is
not a bound on all recovery times. No experiment in this paper is interpreted
as controller-host-loss durability, arbitrary external-side-effect exactly-once
execution, or guaranteed memory safety under undeclared allocation.
